
\section{Computing aggregate error}
\label{appendix:cross_error_aggregate}

%%%%%%%%%%%
%IMPORTANT%

% Version 07/06/2024
% Only one sector (Energy, dw=0.023)
% Update plots and the appendix when we add all sectors
%%%%%%

In this Appendix, we compute the distribution of errors which arise from the combination of the error arising at each step of our forecasting exercise: the industry-level productivity growth (see Appendix \ref{appendix:validation_zeta}), and the Domar aggregation (see Appendix \ref{appendix:validation_Domar}).

We do this at each horizon $\tau$, that is, we assume that we are given i) industry-level productivity growth at each horizon between $t$ and $t+\tau$, and the error associated with that long-term growth, which is a function of $\tau$, ii) fixed Domar weights for year $t$, and iii) aggregate TFP growth rates between $t$ and $t+\tau$ with a certain error, function of $\tau$.

We use the EU-KLEMS data \citep{bontadiniEUKLEMSINTANProdIndustry2023} as in previous appendices.

\subsection{Set-up}


Errors arise at two steps of the forecast: when we predict industry-level productivity growth from prices and quantities, and when we predict aggregate productivity growth using industry-level productivity growth,

First, let us define $\nu_{j,c,t,\tau}$ as the error between the predicted and  true value of productivity growth for country $c$ and industry $j$, 
\begin{equation}
\nu_{j,c,t,\tau} = \hat{A}^{\text{pred}}_{j,c,t,\tau} -\hat{A}_{j,c,t,\tau}.
\label{eq:nu_appendix_agg_error}
\end{equation}
Second, recall from Eq. \ref{eq:varepsiloncttau} in Appendix \ref{appendix:validation_Domar} that we define the error due the Domar aggregation procedure as
\begin{equation}
\varepsilon_{c,t,\tau}=  \hat{A}_{c,t,\tau}^\text{pred} - \hat{A}_{c,t,\tau},
\label{eq:varepsu_appendix_agg_error}
\end{equation}
where the predictions are made using the \textit{true} values of industry-level productivity growth, that is
\begin{equation}
\hat{A}_{c,t,\tau}^\text{pred} = \sum_{j=1}^N \lambda_{j,c,t} \hat{A}_{j,c,t,\tau}.
\label{eq:Acttau_pred_agg_appendix}
\end{equation}
Now, we predict country-level productivity growth as the Domar-weighted average of the \textit{predicted} industry-level productivity growth,
\begin{equation}
 \tilde{\hat{A}}_{c,t,\tau} = \sum_{j=1}^N \lambda_{j,c,t} \hat{A}^{\text{pred}}_{j,c,t,\tau}.
\end{equation}
Using Eq. \ref{eq:nu_appendix_agg_error}, this becomes
\begin{equation}
 \tilde{\hat{A}}_{c,t,\tau} = \sum_{j=1}^N \lambda_{j,c,t} (\hat{A}_{j,c,t,\tau} + \nu_{j,c,t,\tau}) = \hat{A}_{c,t,\tau}^\text{pred} + \sum_{j=1}^N \lambda_{j,c,t} \nu_{j,c,t,\tau},
\end{equation}
where the last step used Eq. \ref{eq:Acttau_pred_agg_appendix}. Finally, inserting  Eq. \ref{eq:varepsu_appendix_agg_error}, this becomes
\begin{equation}
 \tilde{\hat{A}}_{c,t,\tau} =  \hat{A}_{c,t,\tau} + \rho_{c,t,\tau},
\end{equation}
where
\begin{equation}
\label{eq_agg_combined_error}
\rho_{c,t,\tau} \equiv \varepsilon_{c,t,\tau} + \sum_{j=1}^N \lambda_{j,c,t} \nu_{j,c,t,\tau}.
\end{equation}
The goal is now to determine the distribution of $\rho_{c,t,\tau}$ at each horizon $\tau$.

\subsection{Empirical estimation of the distribution of errors}

\paragraph{Data} We use the KLEMS database to calibrate the distribution of errors $\varepsilon$ and $\nu$. 
%We consider the 4 sectors (\er{To be updated}) included in our analysis so far: Energy, Domestic Heating, Steel, Cement with the following Domar weights: Energy = 0.023, Steel = 0.004, Cement = 0.00038, Domestic heating equipment = 0.0035.
In this draft we consider only the Energy sector, with a Domar weight of 0.023. 

%%% To be updated with the number of sector%%%
\paragraph{Method} We draw independently\footnote{
%
Because industries tend to co-move, assuming independence is not innocuous. In this draft, we have only one industry so this does not matter. When we will have several industries, this is an important assumption.
%
} 1 error from the distribution of $\nu_{j,c,t,\tau}$, multiply it by the appropriate Domar weight, and draw one error from the distribution of $\varepsilon_{c,t,\tau}$.%(\ER{number of industries to be updated})%


\subsection{Extrapolation of the quantiles.} 

Before showing the results, we have to deal with one other issue. The KLEMS database is shorter (25Y) than the horizon we plan to forecast (28Y). We thus have to extrapolate the distribution of errors for the longer horizons. We proceed in two steps: first we approximate the distribution of errors per percentiles of the distribution (from 5\% to 95\% with an increment of 5\%). The increment of 5\% is arbitrary and ensures sufficient precision while reducing the noise. We linearly extrapolate the quantiles of errors to years 25 to 35 using the quantiles from horizon 15 to 24.

We validate numerically the extrapolation method step by showing that the percentiles at 5\% increment are a sufficiently good approximation of the distribution of errors, and by plotting the quantiles of errors from horizon 1 to 35 to check that there is no aberration. 

For ease of manipulation, in the code and equations we compute errors in log points. However, for ease of understanding, we plot the error in percentage point.
\footnote{
To transform log growth ($\hat{A}_{t,\tau}=\ln\frac{A_{t+\tau}}{A_{t}}$) into percentage growth ($\frac{A_{t+\tau}}{A_{t}}-1$), we apply the following transformation $\frac{A_{t+\tau}}{A_{t}}-1=\exp (\hat{A}_{t,\tau})-1$
}

In Fig. \ref{g_quantiles}, we plot the quantiles of errors: using KLEMS data (0-25 year) and extrapolated data (26 to 35 years). Note that the quantiles of error at the industry-level seem large, it is due to the long horizon (35 years) and the extreme percentiles represented. The values of the 10-90 percentiles and 25-75 percentiles at horizon 25 years are the same than in appendix \ref{appendix:validation_zeta}.

\begin{figure}[H]
  \centering
   \caption{Quantiles of error for Domar aggregation}
   \label{g_quantiles}
  \subfigure[Industry-level]{
   \includegraphics[width=0.45\linewidth]{figures/hindcasting_agg_error/df_extended_zeta_error.jpeg} 
   }
  \quad
  \subfigure[Domar aggregation]{
     \includegraphics[width=0.45\linewidth]{figures/hindcasting_agg_error/df_extended_hulten_error.jpeg}
     }
     \end{figure}
     
\paragraph{Validation of the errors}

%Figure \ref{fig_histograms_error_horizon} allows to validate the use of the percentiles of error to approximate the distribution for the first 25 years (1000 draws). We then plot the distribution of errors at horizon 30 years, using extrapolated percentiles.
To validate the use of the percentiles of error to approximate the distribution for the first 25 years, we compare the distribution of errors to the outcome of 1000 draws. The results are satisfying. We then plot the distribution of errors at horizon 30 years, using extrapolated percentiles.


% \begin{figure}[H]
%   \centering
%   \caption{Distribution of errors - using the full distribution of the percentiles}
%     \label{fig_histograms_error_horizon}
%   \subfigure[Industry-level]{
%    \includegraphics[width=0.8\linewidth]{figures/hindcasting_agg_error/histograms_error_industry_horizon_1to35.jpeg} 
%    }
%   \quad
%   \subfigure[Domar aggregation]{
%      \includegraphics[width=0.8\linewidth]{figures/hindcasting_agg_error/histograms_error_industry_horizon_1to35.jpeg}
%      }
%      \end{figure}


\subsection{Aggregate error distribution}

We compute the distribution of aggregate error $\mathcal{E}_{\tau}$ from 10,000 draws from the distribution of $\varepsilon_{t,\tau}$ and the draw of 10,000 error for each of the for sectors considered. Figure \ref{fig_error_combined_zeta_domar} shows the resulting distribution. Firstly, we note that the large uncertainty of industry-productivity growth is reduced by the Domar aggregation. Secondly, our estimate is till unbiased, the mean error at the 30 years horizon is close to 0. The 25-50 percentiles is [-3\%,3\%] (in percentage points).

\begin{figure}[H]
    \centering
    \caption{Distribution of errors per horizon - aggregate productivity}
    \label{fig_error_combined_zeta_domar}
    \includegraphics[width=0.8\linewidth]{figures/hindcasting_agg_error/g_error_combined_zeta_domar.jpeg}
\end{figure}


